\documentclass[11pt]{article}
\usepackage[a4paper,margin=25mm]{geometry}
\usepackage[T1]{fontenc}
\usepackage{lmodern}
\usepackage{amsmath,amssymb,amsthm,mathtools}
\usepackage{booktabs,array,longtable}
\usepackage{microtype}
\usepackage[hidelinks]{hyperref}
\usepackage{fancyhdr}
\usepackage{enumitem}
\pagestyle{fancy}
\fancyhf{}
\lhead{Unknot recognition: implementation report}
\rhead{September 2026}
\cfoot{\thepage}
\setlength{\headheight}{14pt}
\setlength{\parskip}{4pt}
\setlength{\parindent}{0pt}
\setlist{nosep,leftmargin=*}
\newtheorem{theorem}{Theorem}[section]
\newtheorem{proposition}[theorem]{Proposition}
\newtheorem{lemma}[theorem]{Lemma}
\newtheorem{corollary}[theorem]{Corollary}
\theoremstyle{definition}
\newtheorem{definition}[theorem]{Definition}
\newcommand{\F}{\mathbb F_2}
\newcommand{\Q}{\mathbb Q}
\newcommand{\Z}{\mathbb Z}
\newcommand{\Kh}{\widetilde{\operatorname{Kh}}}
\newcommand{\code}[1]{\texttt{#1}}
\title{Unknot recognition:\\an exact baseline and partial hierarchy kernels}
\author{Implementation and mathematical audit}
\date{17 September 2026}
\begin{document}
\maketitle

\begin{abstract}
This archive contains an executable exact unknot recognizer, polynomial-time
kernels relevant to the supplied hierarchy notes, tests, and a complexity audit.
\textbf{It does not establish the requested $n^{O(\log n)}$ guarantee.}
The complete recognizer uses reduced Khovanov homology over $\F$ and has an
exponential worst-case bound. The hierarchy kernels test essential boundary
patterns on known balls, compute integral representatives of a rational
cohomology basis, construct compressed cocycle-dual normal surfaces, and check
numerical progress measures. The missing geometric constructions are identified
explicitly rather than represented by unexplained oracle calls. We give
correctness arguments for the implemented parts and a conditional quantitative
lemma describing what a completed quasipolynomial implementation would have to
prove.
\end{abstract}

\setcounter{tocdepth}{1}
\tableofcontents
\newpage

\section{Outcome, sources, and the boundary of the claim}

The requested algorithm is a total decision procedure for classical knots whose
worst-case running time, measured in the number $n$ of crossings in a supplied
diagram, is $n^{O(\log n)}$. The supplied talk announces this theorem and develops
a hierarchy-based route to it \cite{slides}. Its final flowchart, on PDF page
109, is not reproduced here as though each named box were an implemented
subroutine.

The delivered software has two distinct parts. The \emph{complete recognition
path} is an independently implemented cube-of-resolutions algorithm. The
\emph{hierarchy path} consists of genuinely implemented but partial kernels.
There is no complete hierarchy constructor, compressed hierarchy-cutting
engine, or geometric Cheeger-region processor in this archive. In particular,
no claimed depth bound is passed as a numerical parameter and then treated as a
proved topological fact.

Additional research located Lackenby's July 2026 paper \cite{hierarchies}.
Section 9 explicitly leaves some individual step costs unspecified and proves
an iteration bound $L(g+1)^L$ in terms of two parameters. This report does not
interpret that result as a complete bit-complexity specification, and does not
claim that the announced quasipolynomial theorem is false.

Every recognition result contains
\begin{verbatim}
"quasipolynomial_guarantee": false
\end{verbatim}
This is an essential qualification of the deliverable, not a configurable
warning that can be removed by setting a timeout.

\section{Running the implementation}

The source package needs Python 3.10 or later and only its standard library.
The recorded tests used Python 3.13.5 on Linux. From the archive root:
\begin{verbatim}
python -m unknot_recognition recognize examples/torus_3_5.json --check-d2
python -m unknot_recognition recognize examples/hard_unknot.json
python -m unittest discover -s tests -v
python scripts/experiments.py
\end{verbatim}
The first command returns \code{NONTRIVIAL}; the second returns \code{UNKNOT}.
The experiment script regenerates the recorded measurements, including a
separate exhaustive family check.

The public Python entry point is:
\begin{verbatim}
from unknot_recognition import braid_closure, recognize
knot = braid_closure(3, [1, 2] * 5)
result = recognize(knot, check_d_squared=True)
\end{verbatim}
The default pipeline performs local simplifications and a determinant rejection
before using the complete fallback. Both preprocessors can be disabled. Thus
the fallback can be tested independently of the shortcuts.

\subsection{Decisions, limits, and invalid data}

A valid completed answer has status \code{UNKNOT} or \code{NONTRIVIAL} and a
Boolean \code{is\_unknot}. Optional \code{max\_resolutions} and
\code{max\_basis} limits instead produce \code{UNKNOWN} with
\code{is\_unknot: null} when the fallback cannot proceed within them. With both
limits absent, the mathematical algorithm is total; finite hardware still
imposes ordinary memory and execution limits.

The command-line exit code is zero for either completed verdict, two for invalid
input or an I/O failure, three for a resource-limited unknown result, and four
for a failed internal consistency check. A multi-component link is rejected; it
is not converted into a knot by silently dropping components. A non-spherical
rotation system is rejected; it is not fed to a theorem about classical knots.

The interface reads JSON diagrams or braid closures, not drawings, DT codes,
or Gauss words. Normal-surface input is a separate, explicitly triangulated
manifold. No diagram-to-exterior conversion is supplied in the hierarchy path.

\section{Diagram representation and safe preprocessing}

\subsection{A combinatorial sphere diagram}

A crossing is a four-tuple $(a,b,c,d)$ of edge labels in counterclockwise order.
The opposite pair $(a,c)$ is the under-strand; $(b,d)$ is the over-strand. Every
edge occurs twice in the entire collection. The empty collection denotes one
crossing-free circle.

There are $4n$ half-edges, also called darts. Let $\alpha$ pair the two darts of
each edge, and let $\sigma$ rotate darts around each crossing. Following
$\sigma\alpha$ gives the face cycles. Following $\alpha$ and the opposite-dart
pairing gives the link components. The implementation requires one such
component and
\[
 V-E+F=n-2n+F=2.
\]
The rotation system defines an orientable ribbon surface; capping its faces
and applying this Euler test validates the specified spherical embedding. It
is insufficient merely to test whether the abstract graph has some planar
embedding: that embedding need not be the input diagram.

Arbitrary integer edge labels are normalized. The complexity bounds below use
canonical labels of $O(\log n)$ bits. For unusually long literal labels, the
cost of reading and normalizing their actual bit length must be added.

The braid converter tracks strand endpoints with a disjoint-set structure and
checks that the closure permutation has one cycle. A positive Artin generator
makes the left strand the over-strand. A one-component closure on $m$ strands
requires at least $m-1$ crossings, allowing impossible oversized strand counts
to be rejected before allocating an endpoint array.

\subsection{Reidemeister reductions}

An empty one-edge face gives a Reidemeister I deletion. A two-edge face at two
distinct crossings gives a Reidemeister II deletion only when each side of the
bigon has the same over/under status at its two endpoints. Thus one strand is
over twice and the other is under twice. A clasp with the opposite status
pattern is not deleted.

For a valid deletion, opposite edge labels at each removed crossing are joined.
The remaining diagram is reconstructed and revalidated. Every move strictly
lowers the crossing number, so this preprocessing terminates in at most $n$
moves. Repeated scans and reconstruction take polynomial time. No Reidemeister
III move, crossing-increasing search, or completeness claim is attached to this
simplifier.

The saved trace names the face darts and deleted vertices in each current
normalized diagram. The verifier recomputes legal moves before replaying a
step. A valid trace ending in the empty diagram is an explicit unknot proof;
a trace ending earlier proves only equivalence to the remaining diagram.

\subsection{The determinant filter}

Joining opposite over-half-edges produces Wirtinger arcs. At a crossing, with
over-arc $o$ and under-arcs $u,v$, the Fox coloring equation is
\[
 2x_o-x_u-x_v=0.
\]
These equations form an $n\times n$ integer matrix. The absolute value of a
first minor is the knot determinant \cite{coloring}. The code evaluates a minor
by fraction-free Bareiss elimination, including row pivoting and exact-division
checks. The zero-crossing determinant is defined as one.

The only inference made is
\[
 \det K\ne1\quad\Longrightarrow\quad K\text{ is nontrivial}.
\]
Determinant one never causes an unknot answer. In particular, the $T(3,5)$ test
has determinant one and must pass to the complete fallback. Matrix dimensions
and coefficient bit growth give a polynomial-time rejection filter; it does
not settle the difficult determinant-one cases.

\section{The complete reduced Khovanov computation}

We use the usual cube/Frobenius construction \cite{khovanov,barnatan}; the code
is independently written. Only total homology dimension is needed, so the
orientation-dependent shifts of the homological and quantum gradings are not
computed. The reported degree is explicitly the \emph{unnormalized cube degree},
not a normalized bigraded knot invariant.

\subsection{Resolutions and basis indices}

Choose one marked edge. A state $s\in\{0,1\}^n$ resolves every crossing. At a
crossing $(a,b,c,d)$, the zero smoothing joins $(a,b),(c,d)$ and the one smoothing
joins $(a,d),(b,c)$. A disjoint-set calculation obtains the resulting circles,
each stored as a bitset of original edge labels. The marked circle is placed
first in a deterministic ordering.

Let $c(s)$ be the number of circles. Put $A=\F[x]/(x^2)$ and let the marked
circle carry the one-dimensional subspace $\F x$. The state space is
\[
 C_s=\F x\otimes A^{\otimes(c(s)-1)},\qquad
 \dim C_s=2^{c(s)-1}.
\]
We set $C^k=\bigoplus_{|s|=k}C_s$. A tensor basis is an integer bit mask:
zero means the label $1$, one means $x$, and the marked factor is fixed to $x$.
Offsets locate each state basis inside its chain group.

\subsection{Saddles and the differential}

For every state coordinate changing from zero to one, the local saddle merges
two circles or splits one. Comparing circle-support bitsets identifies the
unchanged factors and the active factors. The classical-planarity validation is
important here: an unexpected one-circle-to-one-circle local change is rejected
as an internal inconsistency.

A merge uses
\[
 m(1\otimes1)=1,\quad m(1\otimes x)=m(x\otimes1)=x,
 \quad m(x\otimes x)=0,
\]
and a split uses
\[
 \Delta(1)=1\otimes x+x\otimes1,\qquad \Delta(x)=x\otimes x.
\]
The other tensor factors are copied with their reindexed circle positions.
These formulas preserve the marked-$x$ subspace. If the marked circle merges,
the product is either $x$ or zero; if it splits, both resulting circles carry
$x$. An unmarked split cannot alter the marked factor.

The differential is the sum of all these increasing cube-edge maps. Over $\F$
there is no sign bookkeeping: the two contributions around every cube face
cancel. More explicitly, associativity, coassociativity, commutativity,
cocommutativity, and the Frobenius compatibility
\[
 (m\otimes1)(1\otimes\Delta)=\Delta m
 =(1\otimes m)(\Delta\otimes1)
\]
identify the two composites. These identities can be checked directly on the
basis $1,x$. Consequently $d^{k+1}d^k=0$.

Each differential column is a Python integer. Bit position $j$ stores its
coefficient in target basis vector $j$. Addition is exclusive-or. The
implementation constructs columns on demand, avoiding an additional full
unreduced tensor complex. An optional diagnostic explicitly evaluates $d^2$ on
every source basis vector; all reported small-example runs used this check.

\subsection{Exact homology ranks}

For a stream of column bitsets, Gaussian elimination stores at most one pivot
column for each highest set bit. If a pivot already exists, it is exclusive-ored
away. Each elimination lowers the highest set bit, so the procedure terminates
and counts an exact basis for the column space. Let $r_k=\operatorname{rank}d^k$,
with $r_{-1}=r_n=0$. Then
\[
 h_k=\dim C^k-r_{k-1}-r_k,
 \qquad h=\sum_k h_k.
\]
Negative dimensions are rejected as an internal error. The code also checks
that the unnormalized reduced complex has Euler characteristic $\pm1$, as it
must for a knot. None of these arithmetic checks alone is claimed to replace
the mathematical construction or its invariance theorem.

\section{Correctness and the actual exponential bound}

\begin{theorem}[Detection criterion used by the implementation]
For a classical knot $K$,
\[
 K\text{ is the unknot}\quad\Longleftrightarrow\quad
 \dim_{\F}\Kh(K;\F)=1.
\]
\end{theorem}
\begin{proof}
The deep external input is the Kronheimer--Mrowka theorem: a knot is the unknot
exactly when its reduced rational Khovanov homology has rank one
\cite{detector}. For a finite free integral chain complex, the universal
coefficient exact sequence implies
\[
 \dim_{\F}H(C\otimes\F)\ \ge\ \dim_{\Q}H(C\otimes\Q).
\]
The reduced rational rank of a knot is nonzero: its Euler characteristic after
evaluating the quantum variable at one is $\pm1$. Therefore an $\F$ rank of one
forces a rational rank of one. Conversely, the unknot's reduced integral
homology consists of a single free generator, so its reduced $\F$ rank is one.
The marked-subcomplex convention differs at most by conventional grading
shifts, irrelevant to total dimension.
\end{proof}

\begin{corollary}
On a valid finite one-component classical diagram, with resource caps absent,
\code{recognize} terminates and returns the correct Boolean answer.
\end{corollary}
\begin{proof}
The local reductions preserve the knot. The determinant filter has only the
sound rejection rule given above. Otherwise the finite resolution complex is
constructed and its homology ranks are computed exactly. The detection theorem
then gives both implications of the final decision. With caps present, an
interrupted construction gives no Boolean answer and is reported as unknown.
\end{proof}

\begin{proposition}[Conservative bit-complexity bound]
For normalized $n$-crossing knot diagrams, the complete implementation has
running time $2^{O(n)}$ and memory $2^{O(n)}$. These are exponential upper bounds,
not quasipolynomial bounds.
\end{proposition}
\begin{proof}
Smoothing one crossing changes the number of components by at most one.
Starting with a single knot component and smoothing all $n$ crossings gives
$c(s)\le n+1$ for every complete state. Hence the total chain dimension $D$
satisfies
\[
 D=\sum_{s\in\{0,1\}^n}2^{c(s)-1}\le 2^n2^n=4^n.
\]
There are $2^n$ states. Circle and saddle data take polynomial work and storage
per state or cube edge. A column has at most $D$ bits. Constructing it and
performing at most $D$ pivot eliminations has a conservative polynomial cost in
$n$ and $D$; ordinary bitset elimination is bounded by $O(D^3)$ bit operations
for a matrix with at most $D$ rows and columns. Storing pivots takes at most
$O(D^2)$ bits, plus the state and saddle data. The optional exhaustive $d^2$
check also has cost polynomial in $n,D$. Substituting $D\le4^n$ gives the result.
The polynomial preprocessing does not change this bound.
\end{proof}

This proof is deliberately a loose upper bound for the actual implementation.
Small timings, compressed integer columns, or a quick determinant rejection on
many inputs do not establish a better worst-case exponent. In particular,
$2^{O(n)}$ is not another notation for $n^{O(\log n)}$.

\section{Polynomial kernel: essential patterns on a known ball}

\subsection{Precise scope and input}

The talk defines an essential boundary pattern using disks meeting the pattern
at most three times: the boundary curve must cut off an empty disk, a disk
containing a single arc, or a disk containing a tripod \cite[pp.~24--25]{slides}.
Our kernel assumes the three-manifold is \emph{already known to be a 3-ball}.
On its boundary sphere, every simple closed curve bounds a disk that can be
pushed slightly into the ball. This makes the essentiality question a finite
embedded-graph problem.

Input consists of cyclic edge orders at cubic vertices and a count of
vertex-free circle components. Loops occur twice in a vertex rotation and
parallel edges are allowed. The rotation surface of every connected graph
component is checked to be spherical. A disconnected pattern does not require
nesting data merely to decide essentiality: a separating curve disjoint from
the pattern already gives a violation.

This routine is neither a ball recognizer nor a certificate that a supplied
pattern really came from a particular knot exterior. In particular, $b_1=0$
by itself is not used as a ball-recognition rule.

\begin{definition}
A \emph{bond} of a connected graph is a nonempty minimal edge cut. Equivalently,
its removal leaves exactly two connected components and every removed edge
joins those components.
\end{definition}

\begin{proposition}[Finite criterion]
The empty pattern and a single vertex-free circle are essential. Any pattern
with at least two connected components is inessential. A connected cubic
spherical graph is essential exactly when it has no bond of size one or two
and every bond of size three isolates a single vertex.
\end{proposition}
\begin{proof}
The empty case is immediate. A simple closed curve meets a single embedded
circle an even number of times; with at most three intersections it either
misses it or meets it twice. At least one complementary disk contains the
empty pattern or one arc. For a disconnected pattern, a suitable boundary
component of a small regular neighborhood separates nonempty parts of the
pattern without intersecting it, giving a violating disk.

Now suppose the pattern graph is connected. In a cellular spherical embedding,
a primal bond is dual to a simple cycle, including a loop for a bridge and a
digon for a two-edge bond. Realize this dual cycle transversely to the primal
edges. A one- or two-edge bond has pattern vertices on both sides, so neither
side is an allowed empty disk, arc, or tripod. A three-edge bond is likewise
violating unless one side is a single cubic vertex with its three outgoing
edges, which is exactly a tripod.

For the converse, let a curve with at most three intersections violate the
condition. If one side contains no pattern vertices, its intersection with
the connected pattern consists of at most one properly embedded arc, or is
empty, which is allowed. Thus both sides contain vertices. The graph edges
whose endpoints lie on opposite sides form an edge cut of size at most three;
repeated crossings of an edge only reduce this number. Every nonempty edge
cut contains a bond. If there are no bonds of size one or two, this forces
exactly three distinct crossing edges. The inside and outside induced
subgraphs must each be connected: otherwise a component would have a smaller
edge cut. Hence these three edges form a bond. If every such bond isolates a
vertex, the corresponding side of the original curve is precisely a tripod,
again contradicting violation. This proves the criterion.
\end{proof}

\subsection{Algorithm, witnesses, and cost}

The implementation enumerates every edge subset of size one, two, or three,
runs graph search after deleting it, and checks the bond condition. A
three-edge bond isolating a vertex is discarded. A remaining bond supplies a
dual cycle: the face IDs on the two sides of each edge are linked in cyclic
order. A separate verifier checks the cut, its vertex side, the nontriviality
condition, and every successive dual-face incidence.

For $E$ edges and $V$ vertices the time bound is
\[
 O\bigl((E+E^2+E^3)(V+E)\bigr)=O(E^4)
\]
for cubic graphs. Auxiliary space is linear in graph size. A disconnected
negative result is justified by the component count, not by a fabricated
geometric disk coordinate list. Positive results use the complete finite search.
The tetrahedral and cubical patterns pass; a triangular prism yields a
nontrivial three-edge cut; a dumbbell yields a bridge witness.

\section{Polynomial kernel: cocycles and compressed normal surfaces}

\subsection{Manifold validation and cohomology}

This kernel accepts an ordinary finite simplicial complex specified by
four-distinct-vertex tetrahedra, not generalized or ideal face-pairing data.
It rejects duplicate tetrahedra, faces with more than two incident tetrahedra,
and disconnected complexes. Coherent tetrahedron orientations are solved by
sign propagation across interior faces.

Each vertex link is checked to be a connected triangulated surface: its edge
incidences are at most two, and each vertex fan is a circle or an interval.
A closed link must have Euler characteristic two. A link with boundary must
have Euler characteristic one and a single boundary cycle. Together with the
surface checks, the classification of compact surfaces identifies these as
spheres and disks. This rejects non-manifold vertex gluings rather than merely
checking face incidences.

Orient every edge from its smaller vertex label to its larger label. A
cochain $c$ is a cocycle precisely when
\[
 c_{ab}+c_{bc}-c_{ac}=0\qquad(a<b<c)
\]
on every triangular face. Choose a spanning tree and require its edge values
to be zero. Every cohomology class has a unique representative in this gauge:
a vertex potential kills arbitrary tree values, and a potential constant along
a connected spanning tree is constant everywhere.

Exact rational row reduction on the remaining variables computes a basis.
Clearing denominators and dividing coordinate gcds gives integral cocycles
representing a $\Q$-basis of $H^1$. This is not advertised as a unimodular
$\Z$-basis. The nullity is $b_1$. Rational row-reduction bit growth is polynomial
in the finite incidence matrix's size; all arithmetic uses exact integers or
fractions.

\subsection{The dual normal vector without disk expansion}

An integral cocycle defines a map to $\mathbb R/\Z$. In a tetrahedron with
vertices $v_0<\cdots<v_3$, choose local vertex heights
\[
 h_0=0,\qquad h_i=c_{v_0v_i}\quad(1\le i\le3)
\]
and extend affinely. The cocycle equations make these local maps agree modulo
integer translation on shared faces. The inverse image of $1/2\pmod\Z$ is a
properly embedded, cooriented normal surface.

Sort the four heights as
$h_{i_0}\le h_{i_1}\le h_{i_2}\le h_{i_3}$.
Between successive heights there are exactly
$h_{i_j}-h_{i_{j-1}}$ half-integer levels. The lowest interval gives triangles
at vertex $i_0$, the middle interval gives quadrilaterals separating the two
lowest vertices from the two highest, and the highest interval gives triangles
at vertex $i_3$. Zero gaps give zero disks. Thus seven binary-encoded counts per
tetrahedron describe the entire surface, regardless of their numerical sizes.

\begin{proposition}
The constructed vector satisfies the normal matching equations and the
quadrilateral constraint. Its number of intersections with a global oriented
edge $e$ is $|c_e|$. Its relative homology class is dual to the cocycle class.
\end{proposition}
\begin{proof}
Only one middle partition occurs in a tetrahedron, proving the quadrilateral
constraint. The face restrictions differ only by integer translation, so their
half-integer level arcs coincide; this proves matching. A linear lift changing
by the integer $c_e$ crosses exactly $|c_e|$ half-integer levels along an edge.
Counting signed crossings of a transverse curve computes evaluation of $c$,
which identifies the Poincare-dual relative class.
\end{proof}

Nonzero dual class ensures that some connected component is nonseparating, but
the program does not find that component. A primitive cohomology class is not
used as a substitute for a connectedness test. Nor does a polynomial bit length
of a normal vector imply a polynomial genus or polynomially many disks.

\subsection{Euler characteristic and a compression test}

The validator computes normal arc counts of each type on each face and
requires equality across interior faces. It also checks common global-edge
intersection counts. If $F$ is the total disk count, $A$ is the total disk-side
arc count, and $A_{\partial}$ counts unpaired boundary arcs, the surface has
\[
 E_S=\frac{A+A_{\partial}}2,
 \qquad V_S=\sum_{\text{global edges }e}|S\cap e|,
 \qquad \chi(S)=V_S-E_S+F.
\]
These are integer linear calculations; no disk is enumerated.

A supplied nine-tetrahedron triangulation of triangle $\times$ circle is a
solid torus. Its computed cocycle basis has one element. Scaling this element
by $2^{400}$ produces coordinates representing $3\cdot2^{400}$ disks with
Euler characteristic $2^{400}$. The coordinate payload uses 1263 bits in the
recorded run. This demonstrates only the compressed construction and Euler
calculation, not compressed cutting, normalization, connected-component
extraction, or the Agol--Hass--Thurston orbit-counting algorithm.

\section{What a quasipolynomial proof would need}

\subsection{A precise numerical lemma}

The talk uses an informal genus-like digit complexity, with a warning that a
pattern-sensitive variant may be intended \cite[pp.~54--65]{slides}. The
numerical helper accepts actual nonnegative integer digits $q_i\le g$ and pads
them to length $L$. Define
\[
 \Phi(q_1,\ldots,q_L)=\sum_{i=1}^L q_i(g+1)^{L-i}.
\]
Base $g+1$ is necessary when the digit $g$ is allowed. With base $g$, for
example, $(1,0)$ and $(0,g)$ have equal value. We make this convention explicit
rather than silently interpreting the slide notation differently.

\begin{lemma}[Rebuild bound]
Suppose a rebuild preserves all digits before position $j$, lowers $q_j$ by at
least one, and replaces later digits by arbitrary values in $[0,g]$. Then
$\Phi$ drops by at least one. There are at most $(g+1)^L$ such rebuild epochs;
if at most $L$ construction steps occur per epoch, the step count is at most
$L(g+1)^L$.
\end{lemma}
\begin{proof}
The leading decrease contributes at least $(g+1)^{L-j}$. The largest possible
tail increase is
\[
 g\sum_{k=0}^{L-j-1}(g+1)^k=(g+1)^{L-j}-1.
\]
Their difference is at least one. Finally $0\le\Phi<(g+1)^L$. Counting at most
$L$ construction steps between decreases gives the result.
\end{proof}

The helper also computes the expression
$-4\chi(S)+|S\cap P|$, the pattern complexity used in
\cite[Definition 3.4]{hierarchies}. It does not claim that this expression is
nonnegative for arbitrary surfaces; exceptional negative-complexity disks
must be handled before a nonnegative-digit argument applies.

\subsection{An explicit conditional complexity statement}

\begin{theorem}[Implementation obligations, not a theorem about this software]
Put $N=\max(2,n)$. Suppose a correctly specified geometric recognizer maintains
hierarchies with digits $q_i\le N^a$, length $L\le b\log_2N+c$, and at most
$N^r$ global restart epochs. Suppose also that every elementary operation,
including updates, normalization, certificate checks, and construction of the
next state, has bit cost at most $N^{d\log_2N+e}$, uniformly over every reachable
state. If the preceding rebuild lemma applies within each restart epoch, then
the recognizer runs in $N^{O(\log N)}$ time.
\end{theorem}
\begin{proof}
The total is bounded by an initialization cost plus
\[
 N^r L(N^a+1)^L N^{d\log_2N+e}.
\]
Taking base-two logarithms gives at most
$(ab+d)(\log_2N)^2+O(\log N+\log\log N)$.
This is the required quasipolynomial form.
\end{proof}

The quantified uniformity over reachable states matters. Polynomial cost in
a compressed state's current size is insufficient unless the size and all
integer bit lengths are themselves bounded throughout the complete run.
Similarly, a hierarchy potential may restart at its maximum every time a
Heegaard structure changes; those global restarts need a separate quantitative
bound. A possible combined proof would place a bounded global rank ahead of
$\Phi$ lexicographically, with explicit bounds on both. No such geometric rank
is constructed here.

\subsection{The geometric gap is not a missing loop counter}

The four speedups on pages 66--74 of the talk are bounded surface complexity,
efficient compressed encoding, multi-surfaces, and generalized Heegaard
splittings with Cheeger regions. None can be replaced just by entering
$L=\lceil\log n\rceil$ into a function. The terminal flowchart requires concrete
implementations of \emph{Build Hierarchical Multi-Surface}, \emph{Cut Along
Surface}, \emph{Simplify Multi-Surface}, \emph{Use Cheeger Region},
\emph{Haken's Lemma}, and \emph{Weakly Reduce} \cite[p.~109]{slides}.

For example, the displayed inequality
\[
 3g(\partial M_i)\le
 \min\{g(H_i),\ g(H)-g(H_i)\}
\]
is only part of the Cheeger-region condition. The restricted Morse function
must have the specified Heegaard form. The integer comparison routine in this
archive says nothing about those geometric hypotheses, and produces no
modified handle structure.

An arbitrary integral cocycle construction also does not establish the talk's
quadratic surface-complexity bound: that construction uses the ambient
embedding in $S^3$, which must be preserved through the hierarchy. Binary
coordinates alone provide neither the embedding nor a compressed normalization
algorithm after a disk compression.

\section{Validation and reproducible observations}

The delivered log records \textbf{58 passing test methods}. These test exact
Bareiss elimination against the permutation formula, bitset rank against
independent dense elimination, input rejection, local-move traces, the homology
complex, mirrors, basepoints, Markov stabilization, and the three-strand braid
relation. Fifty seeded random one-component braid closures test invariance
under the implemented reductions, including exhaustive $d^2$ checks before
reduction. Separate tests cover graph witnesses, invalid manifold links,
normal matching, huge normal coordinates, numerical potentials, and CLI exit
semantics.

An additional experiment checks \textbf{682 two-strand braid words}: every word
of odd length at most nine. In $B_2$, the exponent sum determines the braid, so
these closures reduce independently to $T(2,e)$, with $e$ odd. They are unknots
exactly when $|e|=1$. All computed ranks equal $|e|$, and all verdicts agree
with this independent family classification.

\begin{center}
\small
\begin{tabular}{lrrrrl}
\toprule
Input & $n$ & determinant & reduced rank & raw basis & verdict\\
\midrule
Circle & 0 & 1 & 1 & 1 & unknot\\
Curl & 1 & 1 & 1 & 3 & unknot\\
Trefoil & 3 & 3 & 3 & 15 & nontrivial\\
Figure-eight & 4 & 5 & 5 & 33 & nontrivial\\
$T(2,7)$ & 7 & 7 & 7 & 1095 & nontrivial\\
$T(3,5)$ & 10 & 1 & 7 & 4209 & nontrivial\\
R1/R2-resistant unknot & 8 & 1 & 1 & 1365 & unknot\\
\bottomrule
\end{tabular}
\end{center}
The table gives unreduced-by-preprocessing input sizes and chain basis sizes
for the \emph{reduced homology complex}. Thus ``raw basis'' does not mean the
unreduced Khovanov theory. Complete $d^2$ checks were performed on these raw
complexes. The eight-crossing unknot simplifies to six crossings before the
pipeline invokes homology.

\subsection{Independent proof of the resistant unknot example}

Write $a=\sigma_1$ and $b=\sigma_2$ in $B_3$. The input word is
\[
 \beta=b^{-3}b\,a^2ba=b^{-2}a^2ba.
\]
Using $aba=bab$ gives $a^2ba=abab$. Cyclically moving the final $b$ to the
beginning preserves the closure, so
\[
 \widehat\beta=\widehat{b^{-2}abab}
 =\widehat{b^{-1}aba}=\widehat{ab}.
\]
For the last equality, $b^{-1}ab=aba^{-1}$, another form of the same braid
relation. The closure of $ab$ destabilizes to the closure of $a$ in $B_2$,
then to the one-strand unknot. This proves the expected positive answer
without relying on the computed homology rank or the incomplete R1/R2 trace.

The nontriviality of $T(3,5)$ likewise does not depend on its determinant. The
usual torus decomposition gives its group
$\langle u,v\mid u^3=v^5\rangle$, which surjects onto the nonabelian group
$C_3*C_5$. Hence it cannot be the unknot group $\Z$. These two examples check
both directions of the determinant-one fallback.

\subsection{Limits of the validation}

No external knot package was used to cross-check the implementation, and no
Lean or other machine-checked proof was produced. Tests and arithmetic
consistency checks do not by themselves constitute a formal verification of
all code. The report's mathematical arguments and the cited detection theorem
explain why the specified algorithm is correct; the executable remains ordinary
tested research software. Recorded timing measurements describe one environment
and do not prove asymptotic complexity.

\section{Concrete remaining implementation contracts}

A completed geometric path would need at least the following developments.

\textbf{Representation and initialization.}
Construct a marked knot exterior and a layered handle structure from the input
diagram. Retain the ambient embedding, peripheral information, boundary-pattern
labels, and provenance of later cuts. Give an explicit finite encoding of
parallel families and prove global bit-length bounds.

\textbf{Compressed topology.}
Implement cutting, component extraction with orientation and relative homology,
normalization, and pattern updates in that encoding. A small normal vector is
not an implementation of any of these operations. In particular, connectedness
must not be inferred from Euler characteristic.

\textbf{Bounded multi-surfaces and simplification.}
Construct the correct independent geometric surfaces with the required
complexity bounds. Track their intersections and induced sequential cuts. Lift
violating disks to earlier surfaces, distinguish ordinary compression from
boundary/pattern compression, and preserve later work only in the cases where
that preservation is justified.

\textbf{Heegaard and Cheeger operations.}
Check the full geometric region conditions, construct the required weak
reductions or replacement handle structures, and identify a quantitatively
bounded global progress measure. The logarithmic depth estimate and restart
count must apply to the actual chosen implementation, including exceptional
cases.

\textbf{End-to-end proof.}
Verify the terminal disk/essential-hierarchy criterion with all its manifold
hypotheses and combine operation costs, state sizes, hierarchy rebuilding, and
global restarts. Only then would the conditional theorem in Section 8 yield
the requested running-time claim for executable code.

The included baseline can serve as a small-instance oracle for testing such
future constructions. It cannot supply their missing complexity proof.

\begin{thebibliography}{9}
\bibitem{slides}
Marc Lackenby.
\emph{Unknot recognition in quasi-polynomial time}.
February 2021. User-supplied 109-page slide deck,
\code{quasipolynomial-talk.pdf}.

\bibitem{hierarchies}
Marc Lackenby.
\emph{Incompressible surfaces, hierarchies and unknot recognition}.
arXiv:2607.23350v1, 25 July 2026.
\url{https://arxiv.org/abs/2607.23350}.

\bibitem{detector}
P. B. Kronheimer and T. S. Mrowka.
Khovanov homology is an unknot-detector.
\emph{Publications Math\'ematiques de l'IH\'ES} 113 (2011), 97--208.
\url{https://arxiv.org/abs/1005.4346}.

\bibitem{khovanov}
Mikhail Khovanov.
A categorification of the Jones polynomial.
\emph{Duke Mathematical Journal} 101 (2000), 359--426.
\url{https://arxiv.org/abs/math/9908171}.

\bibitem{barnatan}
Dror Bar-Natan.
On Khovanov's categorification of the Jones polynomial.
\emph{Algebraic \& Geometric Topology} 2 (2002), 337--370.
\url{https://arxiv.org/abs/math/0201043}.

\bibitem{coloring}
Sudipta Kolay.
\emph{Knot Colorings: Coloring and Goeritz matrices}.
arXiv:1910.08044, 2019.
\url{https://arxiv.org/abs/1910.08044}.
\end{thebibliography}
\end{document}
